% 第 9 章：验证与边界局限
\section{验证与边界局限}\label{sec:rel-verify}

\subsection{注册测试}
\begin{center}\footnotesize
\begin{tabular}{@{}L{5.2cm}L{4.0cm}L{5.2cm}@{}}
\toprule
测试目标 & 注册位置 & 覆盖\\
\midrule
\srcpath{test\_reliability\_resolver} & \srcpath{tests/CMakeLists.txt} & MC、FMEA、故障模式 FMEA、解析器、数据策略、信息物理\\
\srcpath{test\_three\_stage\_reliability} & \srcpath{tests/CMakeLists.txt} & 三阶段隔离/重构/修复与指标\\
\srcpath{test\_intelligent\_cyber\_physical\_reliability} & \srcpath{tests/CMakeLists.txt} & CT/PT、保护配合、设备序列、信息共享依赖、在线 DAE、DER-FRT、三级对照\\
\srcpath{test\_resilience\_assessment} & \srcpath{tests/CMakeLists.txt} & 交叉调用 \srcpath{run\_sequential\_mc}/\srcpath{run\_distribution\_fmea}\\
\srcpath{reliability\_workflow\_e2e} & GUI E2E & FLISR 三向敏感性、日历年频率、物理割集、三阶段、在线 DAE 与三级对照\\
\bottomrule
\end{tabular}
\end{center}
\srcpath{reliability\_workflow\_e2e} 的三向检查为方法内敏感性（负荷加压$\Rightarrow$EENS 增、
失自动化$\Rightarrow$EENS 增且保留时延与冻结控制分量、正被动故障率增$\Rightarrow$模式频率与
总 EENS 增），\textbf{不}主张不同后果引擎的 EENS 数值可直接互比。

\subsection{边界与局限（如实标注）}
\begin{itemize}
  \item \textbf{网络后果口径}：MC/FMEA 在纯交流上为
        \srcpath{"ac-only-dcopf"}；混合系统为
        \srcpath{"hybrid-acdc-network-lp"}，此时
        \fld{validity.dc\_load\_curtailment\_included}=true、\fld{vsc\_dc\_power\_flow\_modelled}=true。
        两种口径的 \fld{ac\_voltage\_reactive\_feasibility\_certified}=false；
  \item \textbf{重要度语义}：\fld{critical\_components.importance} 是共现/归因秩（多故障态整体切负荷
        归给每个停运元件，可相加超 100\%），\textbf{非} Birnbaum 边际重要度 $\partial\mathrm{EENS}/\partial U_c$；
  \item \textbf{收敛判据}：CoV 仅对 $\bar S>0$ 有定义，零损失运行跑满迭代且记录零 CoV 不代表收敛；
  \item \textbf{尾部风险}：非序贯 VaR/CVaR 为 iid 小时近似，低估天气/修复/储能持续性引致的离散度；
  \item \textbf{信息物理分层}：L1 标量接口矩阵保留为快速筛选；L2 对照入口执行共享通信拓扑、QoS、
        共因环境、物理供电与电池自治、保护/FRT 轨迹和互斥联合事件树，保护误动进入总 EENS/LOLE/LOLF；
  \item \textbf{三阶段范围}：耦合 LinDistFlow（线性化损耗），LCC 与多端口能量路由器被拒
        （清 \fld{ok}）；VSC/DC-DC 恒效率（无固定/二次损耗）；发电机/变压器/换流器/开关故障为选入项；
  \item \textbf{在线保护边界}：在线入口执行 CT/PT、定/反时限、方向、mho/四边形距离、差动、主后备、
        断路器失灵、重合器/熔断器/分段器和 DER-FRT 闭环；其网络为正序、单相故障输入，不认证三相 EMT 或行波保护。
\end{itemize}

\subsection{独立方程 oracle：参数解析器交叉验证}\label{sec:rel-resolver-xref}
注册测试 \srcpath{reliability\_resolver\_cross\_validation} 把生产可靠性参数解析器
\srcpath{src/reliability/reliability\_assessment.cpp:resolve\_reliability\_params} 的规范化输出与一个
不链接 hacdcpf 的独立 Python oracle 对照。证据生成器
\srcpath{tools/reliability\_validation/validate\_reliability\_xref.cpp:emit\_case} 在覆盖全部转换分支的
八个确定性故障模式输入上运行解析器（$\lambda$+MTTR 的运行/日历口径、传统 MTBF 的两种约定、显式 MTTF、
含与不含修复时间的 FOR、以及主动按需），oracle \srcpath{tools/reliability\_validation/run\_cross\_validation.py}
依据 Billinton \& Allan 交替更新闭式独立重算 $U=\lambda/(\lambda+\mu)$（$\mu=H/r$）、MTBF/MTTF$\to\lambda$、
$\lambda=f/((1-f)r)\,H$、日历基改正 $\lambda=f_{cal}/(1-f_{cal}r/H)$ 与主动等效 $\lambda_{act}=\nu_d p_d$，
并逐字段核对解析器输出。

复现命令为
\begin{verbatim}
ctest --test-dir build/macos-release \
  -R '^reliability_resolver_cross_validation$' --output-on-failure
\end{verbatim}

准入阈值为 $10^{-9}$；本轮八个算例六类参数（$\lambda$、修复时间、不可用度、MTTF、日历频率、主动等效频率）
的最大误差为 $0$（独立重算与生产解析器在同一 IEEE~754 双精度下逐位一致），负控制（人为扰动 FOR 分支
$\lambda$ 的 $1\%$）如期触发失败。该 oracle 校核确定性参数解析闭式；系统级 EENS/SAIDI 的 MC 与 FMEA
引擎仍由各自套件验证，本节不以参数解析器一致性冒充系统级外部准确度。

\subsection{证据维护}
本手册为对照 2026-08-17 检出的实现修订；行号引用如与后续代码变更不符，以
\srcpath{include/hacdcpf/reliability/}、\srcpath{include/hacdcpf/analysis/three\_stage\_reliability.hpp}
与注册测试为准。数学模型对照见
\srcpath{docs/theory/reliability\_mathematical\_models\_and\_intelligent\_cyber\_physical\_assessment.md}、
\srcpath{docs/guides/reliability\_calculation\_workflow\_and\_case\_guide.md}、
\srcpath{docs/theory/reliability\_assessment\_models.md}。
